% Fichier complété par tools/complete_missing_corrections.py.
% Source : CIN/CIN-03-Transmetteurs/38_Treuil/38_Treuil.tex
% Statut : rédigé (2 question(s)).
\begin{corrigebox}[Corrigé rédigé]
\CorrigeQuestion{1}
Les deux engrènements donnent
                $\omega_3/\omega_2=-Z_2/Z_{3a}$ puis
                $\omega_4/\omega_3=-Z_{3b}/Z_4$. Le câble ne glisse pas :
                $v_{5/0}=R\omega_4$. Donc
                \[\boxed{v_{5/0}=R\frac{Z_2Z_{3b}}{Z_{3a}Z_4}\,\omega_{2/0}}.\]
\CorrigeQuestion{2}
En posant $k=\omega_4/\omega_2=Z_2Z_{3b}/(Z_{3a}Z_4)$ et
                $k_3=\omega_3/\omega_2=-Z_2/Z_{3a}$,
                \[J_{\rm eq,2}=J_2+J_3k_3^2+J_4k^2+M_5(Rk)^2.\]
                Ramenée à la translation de la charge,
                $M_{\rm eq}=M_5+J_4/R^2+J_3(k_3/(Rk))^2+J_2/(Rk)^2$.
\end{corrigebox}
